\documentclass[12pt,a4paper]{article}

% Stesse impostazioni tipografiche e di pagina del report IVPS.
% Contenuto tratto dal report ESG originale (settembre 2025).
\usepackage[utf8]{inputenc}
\usepackage[T1]{fontenc}
\usepackage[italian]{babel}

\usepackage{amsmath}
\usepackage{booktabs}
\usepackage{graphicx}
\usepackage{float}
\usepackage{geometry}
\usepackage{caption}

\geometry{
    top=2.5cm,
    bottom=2.5cm,
    left=2.7cm,
    right=2.7cm
}

\captionsetup[table]{position=top}
\captionsetup[figure]{position=bottom}
\emergencystretch=2em % Evita un solo sforamento nella lista, senza cambiare lo stile.

\title{
    \textbf{
        ESG e performance nel settore energia:\\
        confronto USA/EU 2013--2024
    }
}
\author{Salvatore Megna}
\date{Settembre 2025}

\begin{document}

\maketitle

\section{Introduzione}

Questo report si occupa di valutare se e come i punteggi ESG (Environmental, Social, Governance) si riflettano sulle performance borsistiche delle principali aziende energetiche in USA ed Europa. L'analisi ha due obiettivi complementari: verificare se i livelli ESG aiutano ad anticipare la performance relativa dei titoli su un orizzonte di 12 mesi, e confrontare, \textit{ex post}, il profilo rischio--rendimento dei portafogli con profilo ESG elevato rispetto a quelli con profilo più basso. Per evitare di confondere l'effetto ESG con l'andamento del comparto energia, i rendimenti sono misurati in forma relativa rispetto a benchmark settoriali coerenti con ciascuna area di mercato.

Per dare significato economico ai risultati, il periodo viene segmentato in tre fasi politico-regolatorie: 2013--2016 (pre-Trump), 2017--2020 (Trump), 2021--2024 (Biden); inoltre distinguiamo sistematicamente USA da Europa, così da cogliere eventuali differenze geografiche e temporali nell'associazione tra sostenibilità e rendimento.

\section{Motivazione economica}

Gli indici ESG condensano informazioni su rischi ambientali, sociali e di governance che spesso non emergono nei bilanci tradizionali. In presenza di punteggi più alti ci si attende: migliore gestione dei rischi regolatori e reputazionali, processi decisionali più solidi, maggiore allineamento alla domanda legata alla transizione energetica. Questi canali possono ridurre la probabilità di incidenti, stabilizzare i flussi di cassa e, nel tempo, abbassare il costo del capitale.

Tre effetti pratici rilevanti per il settore energia:

\begin{enumerate}
\item \textbf{Accesso al capitale:} imprese con profili ESG elevati ottengono condizioni più favorevoli su debito ed equity, anche grazie a strumenti dedicati come green bond e sustainability-linked loan, in cui il pricing è legato a obiettivi di sostenibilità.
\item \textbf{Domanda degli investitori istituzionali:} molti fondi possono acquistare soltanto titoli che rispettano criteri ESG. Un punteggio alto rende il titolo ammissibile a più mandati, genera flussi d'acquisto più ampi e stabili, aumenta la liquidità e tende a ridurre la volatilità.
\item \textbf{Premi di valutazione:} se i rischi percepiti sono inferiori, il costo del capitale richiesto scende. A parità di utili attesi, il valore attuale aumenta: il titolo tende a trattare con multipli più alti e, nel tempo, può mostrare un profilo rischio--rendimento migliore.
\end{enumerate}

Questi benefici non sono automatici. Nel settore energia contano molto il ciclo delle commodity e i tassi: in fasi di forte rialzo del petrolio o di stretta monetaria, i business tradizionali (spesso con punteggi ESG più bassi) possono sovraperformare grazie a cash flow immediati e maggiore potere di prezzo. Inoltre, le metriche ESG differiscono tra provider e possono introdurre rumore di misurazione. Per questo, nel nostro studio misuriamo rendimenti relativi a benchmark settoriali e distinguiamo con cura aree geografiche e periodi politico-regolatori, così da capire quando l'informazione ESG si traduce in vantaggi economici e quando, invece, prevalgono le dinamiche macro-settoriali.

\section{Scelta delle aziende}

\begin{table}[H]
\centering
\caption{Campione delle aziende analizzate (USA/EU): gruppi ESG e ticker Yahoo Finance}
\small
\begin{tabular}{lll}
\toprule
\textbf{Gruppo} & \textbf{Azienda} & \textbf{Ticker (Yahoo Finance)} \\
\midrule
USA -- High ESG & NextEra Energy & NEE \\
USA -- High ESG & Xcel Energy & XEL \\
USA -- High ESG & Duke Energy & DUK \\
USA -- Low ESG & Exxon Mobil & XOM \\
USA -- Low ESG & Chevron & CVX \\
USA -- Low ESG & Occidental Petroleum & OXY \\
EU -- High ESG & Enel & ENEL.MI \\
EU -- High ESG & Iberdrola & IBE.MC \\
EU -- High ESG & E.ON & EOAN.DE \\
EU -- Low ESG & TotalEnergies & TTE.PA \\
EU -- Low ESG & BP & BP.L \\
EU -- Low ESG & ENI & ENI.MI \\
\bottomrule
\end{tabular}
\end{table}

Il campione di aziende analizzato è composto da 12 large cap dell'energia, selezionate per garantire confronto chiaro e dati affidabili in USA ed Europa. In ciascuna area includiamo due blocchi:

\begin{itemize}
\item \textbf{Utility/transizione -- High ESG:} USA (NEE, XEL, DUK) ed Europa (ENEL.MI, IBE.MC, EOAN.DE). Sono società regolamentate o orientate a reti e rinnovabili, con obiettivi climatici pubblici e reporting strutturato.
\item \textbf{Oil \& Gas -- Low ESG:} USA (XOM, CVX, OXY) ed Europa (TTE.PA, BP.L, ENI.MI). Sono grandi gruppi internazionali che operano in tutte le fasi dell'oil \& gas, dall'estrazione alla distribuzione.
\end{itemize}

Perché queste scelte:

\begin{itemize}
\item \textbf{Confrontabilità:} avere, in entrambe le aree, un blocco High e uno Low con modelli di business diversi (utility/transizione vs oil \& gas) permette di leggere i risultati come differenze relative dentro USA ed Europa, senza confonderli con il ciclo del settore.
\item \textbf{Qualità dei dati:} tutte sono large cap incluse in indici principali (S\&P 500, FTSE 100, CAC 40, DAX, FTSE MIB, IBEX 35); la dimensione assicura liquidità e serie storiche complete.
\item \textbf{Copertura informativa:} i punteggi ESG mensili sono disponibili presso Sustainalytics, mentre gli storici di prezzo provengono da Yahoo Finance. Questo consente di costruire un dataset mensile lungo il periodo 2013--2024.
\item \textbf{Lettura economica pulita:} separare utility/transizione da oil \& gas aiuta a capire quando l'eventuale ``segnale ESG'' prevale e quando, invece, contano soprattutto prezzo delle commodity e ciclo macro.
\end{itemize}

In sintesi, il set di aziende massimizza omogeneità e copertura tra aree e gruppi, mantenendo allo stesso tempo la possibilità di confronti chiari su dimensione, liquidità e disponibilità dei dati.

\section{Metodologia}

La nostra analisi procede in quattro passaggi coerenti. Prima costruiamo il dataset longitudinale mensile, integrando prezzi, benchmark settoriali e punteggi ESG. Poi trasformiamo i rendimenti in misure relative al benchmark dell'area. Terzo, testiamo se i livelli ESG anticipano la performance relativa dell'anno successivo. Infine, valutiamo \textit{ex post} il profilo rischio--rendimento dei gruppi (High vs Low) tramite Sharpe annualizzato e confronto statistico.

\subsection{Dataset}

Utilizziamo tre fonti con dati a cadenza mensile: i prezzi dei 12 titoli selezionati (USA ed Europa), i rendimenti dei benchmark settoriali Energy USA (XLE) ed Europa (STOXX Europe 600 Oil \& Gas, codice EXH1.DE), e i punteggi ESG mensili di Sustainalytics per ciascun titolo. Le date sono armonizzate a fine mese.

\subsection{Previsione a 12 mesi}

Per ogni titolo costruiamo il rendimento relativo, sottraendo mese per mese il rendimento del benchmark dell'area geografica di appartenenza. In questo modo isoliamo la componente ``specifica del titolo'' dalla performance del settore energia. L'obiettivo è verificare se i livelli ESG anticipano la performance relativa cumulata dell'anno successivo. Per ogni mese $t$ definiamo il target come la somma dei 12 rendimenti relativi futuri (da $t+1$ a $t+12$). Come variabili esplicative usiamo esclusivamente i punteggi ESG osservati al mese precedente (lag 1: E, S, G e totale). Stimiamo modelli Random Forest distinti per area geografica (USA, Europa) e per finestra temporale coerente con il contesto politico-regolatorio: 2013--2016 (pre-Trump), 2017--2020 (Trump), 2021--2024 (Biden). La validazione è out-of-sample con split temporale: il modello è stato addestrato sui mesi iniziali della finestra temporale analizzata e testato sui mesi finali, con un rapporto tra i mesi di addestramento e i mesi di test di 80/20. Le metriche riportate sono $R^2$ e MAE sul set di test.

\subsection{Profilo rischio--rendimento}

A valle della previsione, valutiamo il profilo rischio--rendimento dei gruppi. Per ogni combinazione periodo $\times$ area $\times$ classe ESG costruiamo una serie equal-weight dei rendimenti relativi mensili e calcoliamo l'indice di Sharpe annualizzato (Formula~\ref{eq:sharpe}). Confrontiamo poi High e Low tramite la differenza di Sharpe e ne stimiamo un intervallo di confidenza al 95\% con bootstrap sui mesi disponibili, riportando anche un p-value a due code. Questo approccio consente di capire se, e quando, i titoli con profilo ESG più elevato mostrano un profilo rischio--rendimento superiore rispetto ai Low, mantenendo costanti area geografica e periodo.

\begin{equation}
\operatorname{Sharpe}_{\mathrm{rel}} =
\frac{\mathrm{E}[R_{\mathrm{titolo}}-R_{\mathrm{settore}}]}{\sigma(R_{\mathrm{titolo}}-R_{\mathrm{settore}})}\sqrt{12}
\label{eq:sharpe}
\end{equation}

\noindent\textit{Formula 1:} indice di Sharpe annualizzato: rapporto tra rendimento in eccesso medio del titolo rispetto al settore e deviazione standard dei rendimenti in eccesso, moltiplicato per $\sqrt{12}$ per annualizzare i dati mensili.

\section{Risultati}

\begin{table}[H]
\centering
\caption{Predittività a 12 mesi calcolata con Random Forest}
\label{tab:rf}
\small
\begin{tabular}{lrrrr}
\toprule
\textbf{Periodo} & \textbf{$R^2$ (EU)} & \textbf{$R^2$ (USA)} & \textbf{MAE (EU)} & \textbf{MAE (USA)} \\
\midrule
2013--2016 & $-1{,}399$ & $-0{,}475$ & $0{,}123$ & $0{,}132$ \\
2017--2020 & $-1{,}813$ & $-0{,}395$ & $0{,}278$ & $0{,}309$ \\
2021--2024 & $-1{,}071$ & $-1{,}387$ & $0{,}163$ & $0{,}201$ \\
\bottomrule
\end{tabular}
\end{table}

Come si vede dalla Tabella~\ref{tab:rf}, i modelli stimati sui punteggi ESG del mese precedente non riescono a prevedere la performance relativa cumulata a 12 mesi: gli $R^2$ out-of-sample risultano negativi in tutti i blocchi temporali e in entrambe le aree, mentre i MAE sono relativamente stabili tra USA ed Europa. Nel nostro campione i livelli ESG non forniscono un segnale predittivo affidabile sui rendimenti relativi a un anno.

\begin{table}[H]
\centering
\caption{Sharpe relativo aggregato per gruppo (equal-weight, annualizzato)}
\label{tab:sharpe}
\small
\begin{tabular}{lrrrrrr}
\toprule
\textbf{Periodo} & \textbf{EU High} & \textbf{EU Low} & $\boldsymbol{\Delta}$\textbf{ EU} & \textbf{USA High} & \textbf{USA Low} & $\boldsymbol{\Delta}$\textbf{ USA} \\
\midrule
2013--2016 & $0{,}152$ & $0{,}112$ & $0{,}040$ & $0{,}431$ & $0{,}094$ & $0{,}337$ \\
2017--2020 & $0{,}894$ & $-0{,}782$ & $1{,}676$ & $0{,}668$ & $0{,}042$ & $0{,}626$ \\
2021--2024 & $-0{,}224$ & $0{,}397$ & $-0{,}622$ & $-0{,}667$ & $0{,}098$ & $-0{,}765$ \\
\bottomrule
\end{tabular}
\end{table}

La Tabella~\ref{tab:sharpe} riassume l'indice di Sharpe annualizzato dei portafogli equal-weight per gruppo (High vs Low ESG) per periodo e area. Nel blocco pre-Trump il vantaggio degli High è limitato e non uniforme. Nel blocco 2017--2020 emerge un profilo rischio--rendimento più favorevole agli High, soprattutto in Europa. Nel blocco 2021--2024 la gerarchia si inverte: i Low ESG ottengono Sharpe mediamente più alti in entrambe le aree.

\begin{table}[H]
\centering
\caption{Robustezza statistica bootstrap del differenziale di Sharpe High--Low ESG}
\label{tab:bootstrap}
\footnotesize
\setlength{\tabcolsep}{3pt}
\begin{tabular}{lrrrrrr}
\toprule
\textbf{Periodo} & $\boldsymbol{\Delta}$\textbf{ EU} & \textbf{IC95\% EU} & \textbf{p-value EU} & $\boldsymbol{\Delta}$\textbf{ USA} & \textbf{IC95\% USA} & \textbf{p-value USA} \\
\midrule
2013--2016 & $0{,}040$ & $[-0{,}918;\,0{,}989]$ & $0{,}930$ & $0{,}337$ & $[-0{,}612;\,1{,}387]$ & $0{,}588$ \\
2017--2020 & $1{,}676$ & $[0{,}371;\,3{,}223]$ & $0{,}530$ & $0{,}626$ & $[-0{,}807;\,2{,}447]$ & $0{,}576$ \\
2021--2024 & $-0{,}622$ & $[-2{,}101;\,0{,}848]$ & $0{,}556$ & $-0{,}765$ & $[-2{,}039;\,0{,}729]$ & $0{,}513$ \\
\bottomrule
\end{tabular}
\end{table}

La Tabella~\ref{tab:bootstrap} riporta i differenziali di Sharpe High--Low con intervalli di confidenza al 95\% e p-value da bootstrap. Il segnale pro-High nel 2017--2020 è il più coerente (soprattutto in Europa). Negli altri periodi, gli intervalli di confidenza comprendono lo zero: ciò implica che, data l'incertezza statistica, il differenziale di Sharpe potrebbe essere positivo, negativo o nullo; in termini formali, non possiamo escludere che $\Delta\mathrm{Sharpe}=0$ ($p>0{,}05$), quindi l'evidenza di un vantaggio stabile è debole.

\section{Conclusioni}

Nel complesso, i pattern osservati non vanno interpretati come un ``effetto presidente'', ma come il riflesso del ciclo macro-settoriale. Nel 2021--2022, per esempio, l'impennata di petrolio e gas ha sostenuto in modo eccezionale i titoli oil \& gas (prevalentemente Low ESG), mentre i rialzi dei tassi e le strozzature nelle catene di fornitura hanno pesato su molte società della transizione. In fasi meno dominate dalle commodity, il vantaggio dei profili ESG più elevati può riemergere o attenuarsi. Da un punto di vista predittivo, i punteggi ESG non anticipano la performance relativa cumulata a 12 mesi; tuttavia, \textit{ex post}, risultano associati a profili rischio--rendimento che si muovono con il contesto, come mostrato dalle tabelle dei risultati. Questo suggerisce cautela nell'uso ``meccanico'' dell'ESG come segnale di timing e rafforza l'idea di integrarlo con variabili macro e di settore (prezzi delle commodity, tassi) e con i fattori di rischio tradizionali. Va inoltre considerato che il nostro campione è concentrato su 12 large cap e su specifiche scelte di benchmark e finestra temporale; un'estensione futura a un universo più ampio, a pesi value-weighted e a diverse specifiche ESG potrebbe chiarire ulteriormente quando e perché l'informazione ESG si traduce in una migliore performance aggiustata per il rischio.

\end{document}
